\subsection{\texorpdfstring{Equicontinuous subsets of \(\mathbf{E}'\)}{Equicontinuous subsets of \textbackslash mathbf\{E\}\textquotesingle{}}}

In this section, E denotes a locally convex space and \(E'\) its dual. Whenever we talk of the polar \(M^{\circ}\) of a set M in E (resp. \(E'\) ), we shall always mean, unless otherwise stated, the polar of M relative to the duality between E and \(E'\) . Recall that if V is a closed convex balanced neighbourhood of 0 in E, we have \(V^{\circ\circ}=V\) (II, p.~45, cor. 3).

PROPOSITION 7. --- Let M be a subset of E'. The following conditions are equivalent :

\begin{enumerate}
\def\labelenumi{(\roman{enumi})}
\item
  M is equicontinuous;
\item
  \(\mathbf{M}\) is contained in the polar of a neighbourhood of 0 in E;
\item
  the polar of \(\mathbf{M}\) is a neighbourhood of 0 in E.
\end{enumerate}

If \(\mathbf{M}\) is equicontinuous, there exists a convex balanced neighbourhood \(V\) of 0 such that \(|u(x)| \leqslant 1\) for all \(x \in V\) and all \(u \in M\) ; then we have that \(\mathbf{M} \subset V^{\circ}\) and (i) implies (ii). With the same notations, if \(\mathbf{M} \subset V^{\circ}\) then \(\mathbf{V} \subset V^{\circ \circ} \subset M^{\circ}\) and (ii) implies (iii). Finally, if \(\mathbf{M}^{\circ}\) contains a convex balanced neighbourhood \(V\) of 0, then \(\mathbf{M} \subset M^{\circ \circ} \subset V^{\circ}\) and the relations \(x \in \varepsilon V\) , \(u \in M\) imply \(|u(x)| \leqslant \varepsilon\) for all \(\varepsilon > 0\) , which proves that (iii) implies (i).

We remark that every \(x \in E\) defines a mapping \(j(x): u \mapsto u(x)\) from \(E'\) into K. Hence we can talk of the S-topology on E, where S is a family of subsets of \(E'\) : this is the inverse image under j of the S-topology on \(K^{E'}\) . We verify immediately that if S is a convex bornology on \(E'\) , then the polars of sets of S form a fundamental system of neighbourhoods of 0 for the S-topology on E. This is so, in particular, when S is the family of equicontinuous subsets of \(E'\) and prop. 7 implies:

COROLLARY 1. --- The topology of E is identical with the topology of uniform convergence on equicontinuous subsets of E'.

More generally, let F be a locally convex space; every \(u \in \mathcal{L}(\mathrm{E}; \mathrm{F})\) defines a map \(j(u):(x, f) \mapsto f(u(x))\) from \(E \times F'\) into K (i.e.~into R or C). This enables us to define, on the space \(\mathcal{L}(\mathrm{E}; \mathrm{F})\) , the topology of uniform convergence on a set of subsets of \(E \times F'\) . In particular:

COROLLARY 2. --- Let S be a family of bounded subsets of E. The S-topology on \(\mathcal{L}(\mathrm{E};\mathrm{F})\) is the topology of uniform convergence on sets of the form \(\mathbf{A}\times\mathbf{B}\subset\mathbf{E}\times\mathbf{F}'\) , where A is in S, and B belongs to the family of equicontinuous subsets of F'.

For every \(u \in \mathcal{L}(\mathrm{E}; \mathrm{F})\) , every \(A \in S\) and every closed convex balanced neighbourhood V of 0 in F, the relation \(u(\mathrm{A}) \subset \mathrm{V}\) is equivalent to \(\ll j(u) (\mathrm{A} \times \mathrm{V}^{\circ})\) is contained in the unit ball of K \(\gg\).

PROPOSITION 8. --- Let H be a family of linear mappings from E into a locally convex space F. For H to be equicontinuous, it is necessary and sufficient that for every equicontinuous subset X in the dual F' of F, the set of linear forms \(f \circ u\) , for \(f \in X\) and \(u \in H\) , is equicontinuous.

It is obvious that the condition is necessary. Suppose it is verified, and let V be a closed convex balanced neighbourhood of 0 in F. Since \(V^{\circ}\) is equicontinuous, there exists a neighbourhood W of 0 in E such that \(\left|f(u(x))\right| \leqslant 1\) for all \(x \in W\) , \(u \in H\) and \(f \in V^{\circ}\) ; in other words, \(u(W) \subset V^{\circ\circ} = V\) for all \(u \in H\) , hence H is equicontinuous.

